\documentclass[11pt,a4paper]{article}
\usepackage[T1]{fontenc}
\usepackage{amsmath}
\usepackage{amssymb}
\usepackage[table]{xcolor}
\usepackage{array}
\usepackage{tabularx}
\usepackage{multirow}
\usepackage{geometry}
\usepackage{enumitem}
\usepackage{parskip}
\usepackage{titlesec}
\usepackage{tocloft}
\usepackage{needspace}
\usepackage[most]{tcolorbox}

\geometry{margin=2.2cm}

\titleformat{\section}{\large\bfseries}{}{0em}{}
\setlength{\cftbeforesecskip}{6pt}

% Question heading that refuses to be orphaned at the foot of a page.
% Reserves 8 lines so the heading stays with the Mark Scheme table beneath it.
\newcommand{\qhead}[2][]{\needspace{8\baselineskip}\section[#1]{#2}}

% Operators and shorthands used across 4037
\DeclareMathOperator{\cosec}{cosec}
\newcommand{\dd}{\mathrm{d}}
\newcommand{\dydx}{\dfrac{\dd y}{\dd x}}
\newcommand{\ncr}[2]{{}^{#1}\mathrm{C}_{#2}}
\newcommand{\npr}[2]{{}^{#1}\mathrm{P}_{#2}}

% Pastel colours, cycled across the questions (q1..q9)
\definecolor{q1color}{HTML}{D6EAF8} % pastel blue
\definecolor{q2color}{HTML}{D5F5E3} % pastel green
\definecolor{q3color}{HTML}{FDEBD0} % pastel orange
\definecolor{q4color}{HTML}{F9E79F} % pastel yellow
\definecolor{q5color}{HTML}{E8DAEF} % pastel purple
\definecolor{q6color}{HTML}{FADBD8} % pastel red/pink
\definecolor{q7color}{HTML}{D1F2EB} % pastel teal
\definecolor{q8color}{HTML}{FCF3CF} % pastel cream
\definecolor{q9color}{HTML}{EBDEF0} % pastel lavender

% Mark-scheme table (FOUR columns: Question | Answer | Marks | Partial Marks or Guidance)
% Optional argument sets the fourth heading; default is Partial Marks.
\newenvironment{mstab}[2][Partial Marks]{%
  {\color{#2!75!black}\nobreak Mark Scheme}\par\nopagebreak\smallskip
  \arrayrulecolor{#2!70!black}%
  \renewcommand{\arraystretch}{1.35}%
  \tabularx{\textwidth}{%
    |>{\columncolor{#2!12}\raggedright\arraybackslash}p{1.6cm}|%
     >{\columncolor{#2!12}\raggedright\arraybackslash}p{5.2cm}|%
     >{\columncolor{#2!12}\centering\arraybackslash}p{1.2cm}|%
     >{\columncolor{#2!12}\raggedright\arraybackslash}X|}%
  \hline
  \rowcolor{#2!45} Question & Answer & Marks & #1 \\ \hline
}{%
  \endtabularx
  \arrayrulecolor{black}%
  \par\medskip
}

% Full-width banner row for Alternative methods: \altrow{q3color}{Alternative 1}
\newcommand{\altrow}[2]{\multicolumn{4}{|l|}{\cellcolor{#1!30}\textbf{#2}} \\ \hline}

% Exemplar-answer box
\newtcolorbox{ansbox}[1]{
  colback=#1!8,
  colbacktitle=#1!30,
  colframe=#1!70!black,
  coltitle=black,
  fonttitle=\normalfont,
  title=Worked Solution,
  boxrule=0.6pt,
  arc=2mm,
  left=4pt, right=4pt, top=4pt, bottom=4pt,
  breakable
}

\newcommand{\topiclabel}[1]{\hfill\normalfont\small\itshape\textcolor{black!55}{#1}}

% Per-question head: \examq{paper-id}{number}{topics}{marks}
% Same command and rendered line as the question paper.
\newcommand{\examq}[4]{\section[#1 - Q#2]{#1 - Q#2 - #3 - #4}}

\begin{document}

\begin{center}
  {\LARGE \textbf{Cambridge O Level Additional Mathematics 4037/21}}\\[6pt]
  {\large May/June 2025 - Paper 2 (Calculator): Mark Scheme \& Worked Solutions}
\end{center}

\bigskip

%%%%%%%%%%%%%%%%%%%
\examq{4037/21/M/J/25}{1}{Quadratic Functions}{5}
%%%%%%%%%%%%%%%%%%%
\begin{mstab}{q1color}
\multirow{2}{*}{1(a)} & $(x+1)^2 - 4$ oe & \multirow{2}{*}{\textbf{B3}} & \textbf{B2} $y = (x+1)^2 - 4$ oe \newline or \newline \textbf{B1} $y = (x+1)^2 + k$, where $k$ is a constant \\ \cline{2-2}
 & $(-1,\,-4)$ & & \\ \hline
1(b) & Correct curve \newline (upward parabola, minimum at $(-1,\,-4)$ in the 3rd quadrant, cutting the $y$-axis at $-3$ and the $x$-axis at $-3$ and $1$) & \textbf{B2} & \textbf{B1} for correct shape in 4 quadrants with minimum in the 3rd quadrant \newline\newline \textbf{B1} for an attempt at a quadratic graph with intercepts marked at $-3$ on the $y$-axis and $-3$ and $1$ on the $x$-axis \\ \hline
\end{mstab}

\begin{ansbox}{q1color}
(a) Stationary point by completing the square\\[4pt]
Halve the coefficient of $x$ and subtract the square of that half:
\[
y = x^2 + 2x - 3 = (x+1)^2 - 1^2 - 3 = (x+1)^2 - 4
\]
Since $(x+1)^2 \ge 0$ for all real $x$, the smallest value of $y$ occurs when $(x+1)^2 = 0$, that is when $x = -1$, giving $y = -4$.
\[
\text{Stationary point} = \boxed{(-1,\,-4)} \quad \text{(a minimum)}
\]

(b) Sketch of the curve\\[4pt]
Intercepts. Setting $x = 0$ gives $y = -3$, so the curve meets the $y$-axis at $(0,\,-3)$. Setting $y = 0$:
\[
x^2 + 2x - 3 = 0 \;\Rightarrow\; (x+3)(x-1) = 0 \;\Rightarrow\; x = -3 \ \text{ or } \ x = 1
\]
so the curve meets the $x$-axis at $(-3,\,0)$ and $(1,\,0)$.

What must be drawn: a parabola opening upwards (the coefficient of $x^2$ is positive), with its minimum at $(-1,\,-4)$, which lies in the third quadrant. The curve passes through the $x$-axis at $-3$ and at $1$, and through the $y$-axis at $-3$, and these three intercepts must be labelled. Reading from left to right the curve descends through the second quadrant, crosses the $x$-axis at $-3$, turns at $(-1,\,-4)$ in the third quadrant, rises through the fourth quadrant crossing the $y$-axis at $-3$, and crosses the $x$-axis again at $1$ before rising into the first quadrant.
\end{ansbox}

%%%%%%%%%%%%%%%%%%%
\examq{4037/21/M/J/25}{2}{Quadratic Functions}{5}
%%%%%%%%%%%%%%%%%%%
\begin{mstab}{q2color}
\multirow{5}{*}{2} & $(3k-2)^2 - 4(2)(k)\,[*0]$ & \textbf{M1} & Uses $b^2 - 4ac$ \newline $*$ could be any inequality sign or $=$ \\ \cline{2-4}
 & $9k^2 - 20k + 4\,[*0]$ & \textbf{A1} & Correctly simplifies to 3-term quadratic \\ \cline{2-4}
 & $(9k-2)(k-2)$ & \textbf{M1} & \textbf{FT} Factorises or solves \textit{their} 3-term quadratic \\ \cline{2-4}
 & Critical values $\dfrac{2}{9}$, $2$ & \textbf{A1} & \\ \cline{2-4}
 & $k < \dfrac{2}{9}$ \newline $k > 2$ & \textbf{A1} & Mark final answer \\ \hline
\end{mstab}

\begin{ansbox}{q2color}
Set of possible values of $k$\\[4pt]
For $2x^2 + (3k-2)x + k = 0$ the coefficients are $a = 2$, $b = 3k-2$, $c = k$. Real and distinct roots require a strictly positive discriminant:
\[
b^2 - 4ac > 0 \;\Rightarrow\; (3k-2)^2 - 4(2)(k) > 0
\]
Expand the square in full and collect terms:
\[
(3k-2)^2 - 8k = 9k^2 - 12k + 4 - 8k = 9k^2 - 20k + 4 > 0
\]
Factorise the 3-term quadratic. Two numbers multiplying to $9 \times 4 = 36$ and adding to $-20$ are $-18$ and $-2$:
\[
9k^2 - 18k - 2k + 4 = 9k(k-2) - 2(k-2) = (9k-2)(k-2) > 0
\]
The critical values are $k = \dfrac{2}{9}$ and $k = 2$. The quadratic in $k$ has a positive leading coefficient, so its graph is a parabola opening upwards and it is positive outside the two roots.
\[
\boxed{k < \tfrac{2}{9} \quad \text{or} \quad k > 2}
\]
\end{ansbox}

%%%%%%%%%%%%%%%%%%%
\examq{4037/21/M/J/25}{3}{Permutations \& Combinations}{5}
%%%%%%%%%%%%%%%%%%%
\begin{mstab}{q3color}
3(a) & $103\,680$ & \textbf{2} & \textbf{M1} for $k \times 5! \times 4! \times 3!$ where $k = 1,\,2,\,3,\,6$ \\ \hline
3(b) & $270$ & \textbf{3} & \textbf{M2} for \newline $\ncr{5}{1} \times \ncr{4}{2} \times \ncr{3}{1} + \ncr{5}{1} \times \ncr{4}{1} \times \ncr{3}{2} + \ncr{5}{2} \times \ncr{4}{1} \times \ncr{3}{1}$ \newline oe \newline\newline OR \newline\newline \textbf{M1} for at least one correct product \newline\newline If 0 scored then \textbf{SC1} for a \underline{final} answer of $540$ or $264$ \\ \hline
\end{mstab}

\begin{ansbox}{q3color}
(a) Arrangements with each discipline together\\[4pt]
Treat each discipline as a single block. There are 3 blocks, so the blocks may be ordered in $3!$ ways. Within the blocks the 5 runners may be ordered in $5!$ ways, the 4 swimmers in $4!$ ways and the 3 gymnasts in $3!$ ways.
\[
3! \times 5! \times 4! \times 3! = 6 \times 120 \times 24 \times 6
\]
\[
= 6 \times 120 = 720, \quad 720 \times 24 = 17\,280, \quad 17\,280 \times 6 = \boxed{103\,680}
\]

(b) Selections of 4 with at least one from each discipline\\[4pt]
With 4 people chosen from 3 disciplines and at least one from each, the split of numbers must be $2$-$1$-$1$ in some order. There are three cases.
\begin{align*}
\text{2 runners, 1 swimmer, 1 gymnast:} && \ncr{5}{2} \times \ncr{4}{1} \times \ncr{3}{1} &= 10 \times 4 \times 3 = 120\\
\text{1 runner, 2 swimmers, 1 gymnast:} && \ncr{5}{1} \times \ncr{4}{2} \times \ncr{3}{1} &= 5 \times 6 \times 3 = 90\\
\text{1 runner, 1 swimmer, 2 gymnasts:} && \ncr{5}{1} \times \ncr{4}{1} \times \ncr{3}{2} &= 5 \times 4 \times 3 = 60
\end{align*}
The three cases are mutually exclusive, so add them:
\[
120 + 90 + 60 = \boxed{270}
\]
\end{ansbox}

%%%%%%%%%%%%%%%%%%%
\examq{4037/21/M/J/25}{4}{Straight-Line Graphs}{7}
%%%%%%%%%%%%%%%%%%%
\begin{mstab}{q4color}
\multirow{7}{*}{4} & $m_{AB} = \dfrac{-9-3}{6--4}$ oe isw or $-\dfrac{6}{5}$ & \textbf{M1} & \\ \cline{2-4}
 & $[m_{\perp} =] \dfrac{5}{6}$ or $\dfrac{-1}{\textit{their } m_{AB}}$ & \textbf{M1} & \textbf{FT} \textit{their} $m_{AB}$ providing that $\dfrac{\text{difference in } y}{\text{difference in } x}$ attempted \\ \cline{2-4}
 & [Midpoint] $\left(\dfrac{-4+6}{2},\, \dfrac{3-9}{2}\right)$ oe isw or $(1,\,-3)$ & \textbf{M1} & \\ \cline{2-4}
 & Finds equation of perpendicular bisector: \newline $y - \textit{their}\,(-3) = -\dfrac{1}{\textit{their}\left(-\frac{6}{5}\right)}(x - \textit{their}\,1)$ oe & \textbf{M1} & \textbf{FT} \textit{their} perpendicular gradient providing gradient is $\dfrac{-1}{\textit{their } m_{AB}}$ and using \textit{their} midpoint \newline\newline \textit{Their} midpoint must not be any of $A$, $B$, $C$ or $D$ \\ \cline{2-4}
 & $11x - 75 = \dfrac{5}{6}x - \dfrac{23}{6}$ & \textbf{M1} & Eliminates one variable \newline \textbf{dep} on previous \textbf{M} mark \newline $11x - 75 = $ \textit{their} perp bisector equation \\ \cline{2-4}
 & $E(7,\,2)$ & \textbf{A1} & \\ \cline{2-4}
 & [Area $CDE$ =] $40$ & \textbf{A1} & Dependent on all previous marks \\ \hline
\end{mstab}

\begin{ansbox}{q4color}
Area of triangle $CDE$\\[4pt]
Gradient of $AB$, with $A(-4,\,3)$ and $B(6,\,-9)$:
\[
m_{AB} = \frac{-9-3}{6-(-4)} = \frac{-12}{10} = -\frac{6}{5}
\]
Gradient of the perpendicular bisector:
\[
m_{\perp} = \frac{-1}{m_{AB}} = \frac{-1}{-\frac{6}{5}} = \frac{5}{6}
\]
Midpoint of $AB$:
\[
\left(\frac{-4+6}{2},\, \frac{3+(-9)}{2}\right) = \left(\frac{2}{2},\, \frac{-6}{2}\right) = (1,\,-3)
\]
Equation of the perpendicular bisector, through $(1,\,-3)$ with gradient $\dfrac{5}{6}$:
\[
y - (-3) = \frac{5}{6}(x - 1) \;\Rightarrow\; y = \frac{5}{6}x - \frac{5}{6} - 3 \;\Rightarrow\; y = \frac{5}{6}x - \frac{23}{6}
\]
The line $L$ is $y = 11x - 75$. At $E$ the two lines meet, so eliminate $y$:
\[
11x - 75 = \frac{5}{6}x - \frac{23}{6}
\]
Multiply throughout by 6:
\[
66x - 450 = 5x - 23 \;\Rightarrow\; 61x = 427 \;\Rightarrow\; x = 7
\]
and then $y = 11(7) - 75 = 77 - 75 = 2$, so $E(7,\,2)$.

Area of triangle $CDE$ with $C(15,\,10)$, $D(14,\,-1)$, $E(7,\,2)$, taken in that order:
\[
\text{Area} = \tfrac{1}{2}\left| x_C(y_D - y_E) + x_D(y_E - y_C) + x_E(y_C - y_D) \right|
\]
\[
= \tfrac{1}{2}\left| 15(-1-2) + 14(2-10) + 7(10-(-1)) \right|
= \tfrac{1}{2}\left| -45 - 112 + 77 \right|
\]
\[
= \tfrac{1}{2}(80) = \boxed{40}
\]
\end{ansbox}

%%%%%%%%%%%%%%%%%%%
\examq{4037/21/M/J/25}{5}{Differentiation}{4}
%%%%%%%%%%%%%%%%%%%
\begin{mstab}{q5color}
\multirow{3}{*}{5} & $\dydx = 2x\tan\dfrac{x}{2} + \dfrac{x^2}{2}\sec^2\dfrac{x}{2}$ oe & \textbf{2} & \textbf{B1} for $\dfrac{\dd}{\dd x}\left(\tan\dfrac{x}{2}\right) = \dfrac{1}{2}\sec^2\dfrac{x}{2}$ soi \newline\newline or \textbf{M1} \textbf{FT} for a product rule of correct structure with \textit{their} $\dfrac{\dd}{\dd x}\left(\tan\dfrac{x}{2}\right)$ \\ \cline{2-4}
 & $\dfrac{\delta y}{h} = \textit{their} \left.\dydx\right|_{x=\frac{\pi}{3}}$ \newline or \newline $\textit{their} \left.\dydx\right|_{x=\frac{\pi}{3}} \times h$ & \textbf{M1} & Correct small changes relationship \\ \cline{2-4}
 & $1.94h$ & \textbf{A1} & Dependent on \textbf{B2} \\ \hline
\end{mstab}

\begin{ansbox}{q5color}
Approximate change in $y$\\[4pt]
The angle is in radians, so the calculator is in radian mode and the standard derivative $\dfrac{\dd}{\dd x}(\tan x) = \sec^2 x$ applies directly.

Differentiate $y = x^2\tan\dfrac{x}{2}$ by the product rule with $u = x^2$ and $v = \tan\dfrac{x}{2}$.

The inner function $\dfrac{x}{2}$ has derivative $\dfrac{1}{2}$, so by the chain rule
\[
\frac{\dd}{\dd x}\left(\tan\frac{x}{2}\right) = \frac{1}{2}\sec^2\frac{x}{2}
\]
Hence
\[
\dydx = 2x\tan\frac{x}{2} + x^2 \cdot \frac{1}{2}\sec^2\frac{x}{2} = 2x\tan\frac{x}{2} + \frac{x^2}{2}\sec^2\frac{x}{2}
\]
Evaluate at $x = \dfrac{\pi}{3}$, so that $\dfrac{x}{2} = \dfrac{\pi}{6}$, where $\tan\dfrac{\pi}{6} = \dfrac{1}{\sqrt{3}}$ and $\cos\dfrac{\pi}{6} = \dfrac{\sqrt{3}}{2}$, giving $\sec^2\dfrac{\pi}{6} = \dfrac{1}{3/4} = \dfrac{4}{3}$:
\[
\left.\dydx\right|_{x=\frac{\pi}{3}} = 2\left(\frac{\pi}{3}\right)\frac{1}{\sqrt{3}} + \frac{1}{2}\left(\frac{\pi}{3}\right)^2 \cdot \frac{4}{3}
= \frac{2\pi}{3\sqrt{3}} + \frac{2\pi^2}{27}
= \frac{2\sqrt{3}\,\pi}{9} + \frac{2\pi^2}{27}
\]
As $x$ increases from $\dfrac{\pi}{3}$ to $\dfrac{\pi}{3} + h$ the increment is $\delta x = h$, and for small $h$ the small changes relationship gives
\[
\delta y \approx \left.\dydx\right|_{x=\frac{\pi}{3}} \times h
\]
\[
\delta y \approx \left(\frac{2\sqrt{3}\,\pi}{9} + \frac{2\pi^2}{27}\right)h
\]
The bracket is kept exact until the end. Evaluating it in radian mode,
\[
\frac{2\sqrt{3}\,\pi}{9} = 1.20920\ldots, \qquad \frac{2\pi^2}{27} = 0.73109\ldots, \qquad \text{sum} = 1.94029\ldots
\]
\[
\boxed{\delta y \approx 1.94h}
\]
\end{ansbox}

%%%%%%%%%%%%%%%%%%%
\examq{4037/21/M/J/25}{6}{Binomial Theorem}{6}
%%%%%%%%%%%%%%%%%%%
\begin{mstab}{q6color}
\multirow{2}{*}{6(a)} & $(3x+2)^{10}$ oe & \textbf{M2} & \textbf{M1} for $(3x+2)^2$ soi \\ \cline{2-4}
 & $1024 + 15\,360x + 103\,680x^2$ & \textbf{A2} & \textbf{A1} for any two terms correct or \newline $2^{10} + 10(2^9)(3x) + 45(2^8)(3x)^2$ \newline or for 3 correct terms seen but not summed e.g. a list \newline\newline After \textbf{M2 A0}, allow \textbf{SC1} for an answer of \newline $59\,049x^{10} + 393\,660x^9 + 1\,180\,980x^8$ \\ \hline
6(b) & $51\,963\,120$ & \textbf{2} & \textbf{M1} for $\ncr{12}{4}\left(\dfrac{6}{x^2}\right)^{8}\left(\dfrac{x^4}{2}\right)^{4}$ soi \\ \hline
\end{mstab}

\begin{ansbox}{q6color}
(a) First three terms of $(9x^2 + 12x + 4)^5$\\[4pt]
The quadratic is a perfect square:
\[
9x^2 + 12x + 4 = (3x)^2 + 2(3x)(2) + 2^2 = (3x+2)^2
\]
Therefore
\[
(9x^2 + 12x + 4)^5 = \left[(3x+2)^2\right]^5 = (3x+2)^{10}
\]
Expand by the binomial theorem in ascending powers of $x$, taking $a = 2$ and $b = 3x$ with $n = 10$:
\[
(2 + 3x)^{10} = 2^{10} + \binom{10}{1}2^{9}(3x) + \binom{10}{2}2^{8}(3x)^2 + \cdots
\]
Term by term:
\begin{align*}
2^{10} &= 1024\\
\binom{10}{1}2^{9}(3x) &= 10 \times 512 \times 3x = 15\,360x\\
\binom{10}{2}2^{8}(3x)^2 &= 45 \times 256 \times 9x^2 = 103\,680x^2
\end{align*}
\[
\boxed{1024 + 15\,360x + 103\,680x^2}
\]

(b) Term independent of $x$ in $\left(\dfrac{6}{x^2} + \dfrac{x^4}{2}\right)^{12}$\\[4pt]
The general term is
\[
\ncr{12}{r}\left(\frac{6}{x^2}\right)^{12-r}\left(\frac{x^4}{2}\right)^{r}
\]
Collect the powers of $x$:
\[
x^{-2(12-r)} \times x^{4r} = x^{-24 + 2r + 4r} = x^{6r-24}
\]
The term is independent of $x$ when $6r - 24 = 0$, that is $r = 4$. Substituting:
\[
\ncr{12}{4}\left(\frac{6}{x^2}\right)^{8}\left(\frac{x^4}{2}\right)^{4}
= 495 \times \frac{6^8}{x^{16}} \times \frac{x^{16}}{2^4}
= 495 \times \frac{6^8}{16}
\]
\[
= 495 \times \frac{1\,679\,616}{16} = 495 \times 104\,976 = \boxed{51\,963\,120}
\]
\end{ansbox}

%%%%%%%%%%%%%%%%%%%
\examq{4037/21/M/J/25}{7}{Functions \textperiodcentered\ Logarithmic \& Exponential Functions}{8}
%%%%%%%%%%%%%%%%%%%
\begin{mstab}{q7color}
7(a) & Correct pair of graphs with \textbf{both} asymptotes and \textbf{both} intercepts correctly shown \newline (decreasing curve $y = \mathrm{f}(x)$ through $(0,\,5)$ with asymptote $y = 3$; its reflection $y = \mathrm{f}^{-1}(x)$ through $(5,\,0)$ with asymptote $x = 3$; line $y = x$) & \textbf{4} & \textbf{B1} for correct shape for f; must have positive $y$-intercept and appear to tend to an asymptote in the 1st quadrant \newline\newline \textbf{B1} for $(0,\,5)$ indicated; must have attempted correct shape for f \newline\newline \textbf{B1} for asymptote at $y = 3$; must have attempted correct shape for f \newline\newline \textbf{B1} for a correct reflection of \textit{their} f in the line $y = x$ \newline\newline $y = x$ does not have to be drawn \newline\newline \textbf{Maximum of B3 if not fully correct} \\ \hline
\multirow{3}{*}{7(b)} & $\mathrm{e}^{y} = \dfrac{3}{2-x} - 2$ oe \newline OR \newline $\mathrm{e}^{x} = \dfrac{3}{2-y} - 2$ oe and variables swapped at some point & \textbf{M1} & Rearranges to $\mathrm{e}^{y} = \ldots$ \newline e.g. $\mathrm{e}^{y} = \dfrac{2x-1}{2-x}$ or $\mathrm{e}^{y} = \dfrac{1-2x}{x-2}$ or $\mathrm{e}^{y} = \dfrac{-3}{x-2} - 2$ \\ \cline{2-4}
 & $\ln\left(\dfrac{3}{2-x} - 2\right)$ oe, isw & \textbf{A1} & \\ \cline{2-4}
 & $1 \le x < 2$ & \textbf{B2} & \textbf{B1} $x \ge 1$ oe \newline OR \newline $x < 2$ oe \newline OR \newline $1 < x \le 2$ \\ \hline
\end{mstab}

\begin{ansbox}{q7color}
(a) Sketch of $y = \mathrm{f}(x)$ and $y = \mathrm{f}^{-1}(x)$\\[4pt]
For $\mathrm{f}(x) = 2\mathrm{e}^{-x} + 3$, $x \in \mathbb{R}$:
\[
\mathrm{f}(0) = 2\mathrm{e}^{0} + 3 = 2 + 3 = 5
\]
so the $y$-intercept is $(0,\,5)$. As $x \to \infty$, $\mathrm{e}^{-x} \to 0$, so $\mathrm{f}(x) \to 3$ from above: the horizontal asymptote is $y = 3$. As $x \to -\infty$, $\mathrm{e}^{-x} \to \infty$, so $\mathrm{f}(x) \to \infty$. The function is decreasing throughout, and since $\mathrm{f}(x) > 3$ for all $x$ the curve never meets the $x$-axis.

What must be drawn for $y = \mathrm{f}(x)$: a single decreasing curve, steep in the second quadrant, crossing the $y$-axis at the labelled point $(0,\,5)$, then flattening in the first quadrant and tending to the labelled horizontal asymptote $y = 3$ without touching it.

The graph of $y = \mathrm{f}^{-1}(x)$ is the reflection of $y = \mathrm{f}(x)$ in the line $y = x$. Reflecting swaps the coordinates of every point and swaps the roles of the axes, so:
\begin{itemize}[nosep]
\item the intercept $(0,\,5)$ becomes the $x$-intercept $(5,\,0)$,
\item the horizontal asymptote $y = 3$ becomes the vertical asymptote $x = 3$,
\item the domain of $\mathrm{f}^{-1}$ is $x > 3$.
\end{itemize}
What must be drawn for $y = \mathrm{f}^{-1}(x)$: a decreasing curve lying to the right of the labelled vertical asymptote $x = 3$, falling steeply from large positive $y$ near $x = 3$, crossing the $x$-axis at the labelled point $(5,\,0)$, and continuing downwards to the right. The line $y = x$ may be drawn to show the reflection but is not required.

(b) $\mathrm{g}^{-1}(x)$ and its domain\\[4pt]
Write $y = \mathrm{g}(x)$ and make $x$ the subject:
\[
y = 2 - \frac{3}{\mathrm{e}^{x} + 2} \;\Rightarrow\; \frac{3}{\mathrm{e}^{x}+2} = 2 - y \;\Rightarrow\; \mathrm{e}^{x} + 2 = \frac{3}{2-y}
\]
\[
\mathrm{e}^{x} = \frac{3}{2-y} - 2 \;\Rightarrow\; x = \ln\left(\frac{3}{2-y} - 2\right)
\]
Swapping the variables,
\[
\boxed{\mathrm{g}^{-1}(x) = \ln\left(\frac{3}{2-x} - 2\right)}
\]
The domain of $\mathrm{g}^{-1}$ is the range of $\mathrm{g}$. The domain of $\mathrm{g}$ is $x \ge 0$. At $x = 0$:
\[
\mathrm{g}(0) = 2 - \frac{3}{\mathrm{e}^{0}+2} = 2 - \frac{3}{3} = 1
\]
As $x \to \infty$, $\mathrm{e}^{x} + 2 \to \infty$, so $\dfrac{3}{\mathrm{e}^{x}+2} \to 0$ and $\mathrm{g}(x) \to 2$ from below without ever reaching it. Since $\mathrm{g}$ is increasing on $x \ge 0$, its range is $1 \le \mathrm{g}(x) < 2$.
\[
\text{Domain of } \mathrm{g}^{-1}: \quad \boxed{1 \le x < 2}
\]
\end{ansbox}

%%%%%%%%%%%%%%%%%%%
\examq{4037/21/M/J/25}{8}{Trigonometry}{9}
%%%%%%%%%%%%%%%%%%%
\begin{mstab}{q8color}
\multirow{2}{*}{8(a)} & $\cos x = 0$ \quad or \quad $\tan x = \dfrac{3}{2}$ oe & \textbf{M1} & \\ \cline{2-4}
 & $[x =]\ \dfrac{\pi}{2}$ or $1.57[07\ldots]$ \newline $0.983$ or $0.9827[93\ldots]$ & \textbf{A2} & and no extra solutions in range \newline \textbf{A1} for either ignoring extras \\ \hline
\end{mstab}

\begin{mstab}{q8color}
\multirow{5}{*}{8(b)} & $\dfrac{2}{\sin(2\theta+1)} - 12\sin(2\theta+1) = 5$ \ or \newline $2\cosec(2\theta+1) - \dfrac{12}{\cosec(2\theta+1)} = 5$ & \textbf{M1} & Writes in terms of $\cosec(2\theta+1)$ or $\sin(2\theta+1)$ only \\ \cline{2-4}
 & $12\sin^2(2\theta+1) + 5\sin(2\theta+1) - 2\,[= 0]$ \newline or \newline $2\cosec^2(2\theta+1) - 5\cosec(2\theta+1) - 12\,[= 0]$ oe & \textbf{A1} & \\ \cline{2-4}
 & $(4\sin(2\theta+1) - 1)(3\sin(2\theta+1) + 2)\,[= 0]$ \newline or \newline $(2\cosec(2\theta+1) + 3)(\cosec(2\theta+1) - 4)\,[= 0]$ & \textbf{M1} & Factorises or solves \textit{their} 3-term quadratic in $\sin(2\theta+1)$ or $\cosec(2\theta+1)$ \newline \textbf{dep} on previous \textbf{M1} \\ \cline{2-4}
 & $2\theta+1 = 0.253$ or $0.2526[80\ldots]$ \newline OR \newline $2\theta+1 = -0.73[0]$ or $-0.7297[27\ldots]$ \newline OR \newline $2\theta+1 = \sin^{-1}\left(\textit{their}\ \dfrac{1}{4}\right)$ \newline OR \newline $2\theta+1 = \sin^{-1}\left(\textit{their}\ -\dfrac{2}{3}\right)$ & \textbf{M1} & \textbf{dep} on previous \textbf{M1} \newline \textbf{FT} providing \newline $-1 \le \textit{their}\ \dfrac{1}{4} \le 1$ and/or $-1 \le \textit{their}\ -\dfrac{2}{3} \le 1$ \\ \cline{2-4}
 & $-0.374$ \ or \ $-0.3736[59\ldots]$ \newline $0.944$ \ or \ $0.9444[56\ldots]$ \newline $-0.865$ \ or \ $-0.8648[63\ldots]$ \newline $1.44$ \ or \ $1.435[66\ldots]$ & \textbf{A2} & and no extra solutions in range \newline\newline \textbf{A1} for two correct solutions ignoring extras \\ \hline
\end{mstab}

\begin{ansbox}{q8color}
(a) Solving $(2 - 3\cot x)\cos x = 0$ for $0 < x \le \dfrac{\pi}{2}$\\[4pt]
The interval is given in radians, so the calculator is in radian mode and answers are given to 3 significant figures.

A product is zero when one factor is zero, so
\[
\cos x = 0 \qquad \text{or} \qquad 2 - 3\cot x = 0
\]
From the first factor, in the interval $0 < x \le \dfrac{\pi}{2}$:
\[
\cos x = 0 \;\Rightarrow\; x = \frac{\pi}{2} \ (= 1.5707\ldots)
\]
From the second factor:
\[
3\cot x = 2 \;\Rightarrow\; \cot x = \frac{2}{3} \;\Rightarrow\; \tan x = \frac{3}{2}
\]
\[
x = \tan^{-1}\!\left(\frac{3}{2}\right) = 0.98279\ldots
\]
which lies in the interval. There are no further solutions in $0 < x \le \dfrac{\pi}{2}$.
\[
\boxed{x = 0.983 \quad \text{or} \quad x = \frac{\pi}{2} \ (1.57)}
\]

(b) Solving $2\cosec(2\theta+1) - 12\sin(2\theta+1) = 5$ for $-1 \le \theta \le 2$\\[4pt]
The interval is $-1 \le \theta \le 2$ with $\theta$ in radians, so the calculator is in radian mode.

Let $u = 2\theta + 1$ and $s = \sin u$. Since $\cosec u = \dfrac{1}{\sin u}$, the equation becomes
\[
\frac{2}{s} - 12s = 5
\]
Multiply throughout by $s$ (note $s \ne 0$) and collect terms:
\[
2 - 12s^2 = 5s \;\Rightarrow\; 12s^2 + 5s - 2 = 0
\]
Factorise. Two numbers multiplying to $12 \times (-2) = -24$ and adding to $5$ are $8$ and $-3$:
\[
12s^2 + 8s - 3s - 2 = 4s(3s+2) - (3s+2) = (4s-1)(3s+2) = 0
\]
\[
s = \frac{1}{4} \qquad \text{or} \qquad s = -\frac{2}{3}
\]
Change the interval for $u$. If $-1 \le \theta \le 2$ then $-1 \le u \le 5$.

Case $\sin u = \dfrac{1}{4}$. The principal value is $u = \sin^{-1}\dfrac{1}{4} = 0.25268\ldots$, and the second solution in the interval comes from $u = \pi - 0.25268 = 2.88891\ldots$ Both lie in $[-1,\,5]$. Then $\theta = \dfrac{u-1}{2}$:
\[
\theta = \frac{0.25268 - 1}{2} = -0.37365\ldots, \qquad \theta = \frac{2.88891 - 1}{2} = 0.94445\ldots
\]

Case $\sin u = -\dfrac{2}{3}$. The principal value is $u = \sin^{-1}\!\left(-\dfrac{2}{3}\right) = -0.72972\ldots$, and the second solution in the interval comes from $u = \pi + 0.72972 = 3.87132\ldots$ Both lie in $[-1,\,5]$. Then
\[
\theta = \frac{-0.72972 - 1}{2} = -0.86486\ldots, \qquad \theta = \frac{3.87132 - 1}{2} = 1.43566\ldots
\]
Any further values of $u$ differ by $2\pi$ and fall outside $[-1,\,5]$, so there are exactly four solutions.
\[
\boxed{\theta = -0.865,\ -0.374,\ 0.944,\ 1.44}
\]
\end{ansbox}

%%%%%%%%%%%%%%%%%%%
\examq{4037/21/M/J/25}{9}{Circular Measure \textperiodcentered\ Differentiation}{12}
%%%%%%%%%%%%%%%%%%%
\begin{mstab}{q9color}
9(a) & $135\pi$ isw or $424$ or $424.1$ to $424.2$ & \textbf{2} & \textbf{M1} for $\dfrac{1}{2} \times 15^2 \times \dfrac{6\pi}{5}$ oe \\ \hline
9(b) & $18\pi$ or $56.5$ to $56.6$ & \textbf{2} & \textbf{M1} for $15 \times \dfrac{6\pi}{5}$ oe \\ \hline
9(c) & radius $= 9$, height $= 12$ & \textbf{2} & \textbf{B1} for radius $= 9$ \newline OR \newline \textbf{M1} for a correct Pythagoras' statement \newline e.g. \textit{their} $r^2 + \text{height}^2 = 15^2$ \\ \hline
9(d)(i) & $r = \dfrac{3}{4}h$ oe & \textbf{B1} & \\ \hline
\end{mstab}

\begin{mstab}{q9color}
\multirow{5}{*}{9(d)(ii)} & $V = \dfrac{1}{3}\pi\left(\dfrac{3}{4}h\right)^{2} \times h$ oe & \textbf{B1} & \textbf{FT} \textit{their} $r = kh$ where $k$ is a positive constant \newline e.g. $V = \dfrac{1}{3}\pi\left(\textit{their } kh\right)^{2} \times h$ \\ \cline{2-4}
 & $\dfrac{\dd V}{\dd h} = \dfrac{9\pi h^2}{16}$ oe & \textbf{B1} & \textbf{FT} \textit{their} $k$ where $k$ is a positive constant \newline e.g. $\dfrac{\dd V}{\dd h} = \left(\textit{their } k\right)^{2}\pi \times h^{2}$ \\ \cline{2-4}
 & Relevant, correct chain rule \textbf{including} $\dfrac{\dd h}{\dd t}$ soi \newline $\dfrac{\dd h}{\dd t} = \dfrac{\dd V}{\dd t} \times \dfrac{\dd h}{\dd V}$ oe & \textbf{B1} & e.g. $\dfrac{\dd h}{\dd t} = \dfrac{\dd V}{\dd t} \div \dfrac{\dd V}{\dd h}$ \\ \cline{2-4}
 & $27 \times \dfrac{16}{9\pi \times 16}$ oe & \textbf{M1} & \\ \cline{2-4}
 & $0.955$ or $0.9549[29\ldots]$ & \textbf{A1} & \\ \hline
\end{mstab}

\begin{ansbox}{q9color}
(a) Curved surface area of the cone\\[4pt]
All angles in this question are in radians, so the calculator is set to radian mode throughout and $\pi$ is kept symbolic until the final evaluation.

The curved surface of the cone is exactly the sector $AOB$, of radius $15$ and angle $\dfrac{6\pi}{5}$ radians. Using $A = \dfrac{1}{2}r^2\theta$:
\[
A = \frac{1}{2} \times 15^2 \times \frac{6\pi}{5} = \frac{1}{2} \times 225 \times \frac{6\pi}{5} = \frac{1350\pi}{10} = 135\pi
\]
Keeping the answer exact, $A = 135\pi$; evaluating, $135\pi = 424.115\ldots$, which to 3 significant figures is $424$.
\[
\boxed{135\pi \,\text{cm}^2 = 424 \,\text{cm}^2} \quad (\text{accept } 424.1 \text{ to } 424.2)
\]

(b) Circumference of the circular top\\[4pt]
When the sector is rolled up, the arc $AB$ becomes the circular rim of the cone. Using $s = r\theta$:
\[
s = 15 \times \frac{6\pi}{5} = \frac{90\pi}{5} = 18\pi
\]
Evaluating, $18\pi = 56.548\ldots$, which to 3 significant figures is $56.5$.
\[
\boxed{18\pi \,\text{cm} = 56.5 \,\text{cm}} \quad (\text{accept } 56.5 \text{ to } 56.6)
\]

(c) Radius of the top and perpendicular height\\[4pt]
The rim of the cone is a circle of circumference $18\pi$, so
\[
2\pi r = 18\pi \;\Rightarrow\; \boxed{r = 9 \,\text{cm}}
\]
The slant height of the cone is the original radius, $15$. By Pythagoras' theorem on the right-angled triangle formed by the radius, the height and the slant height:
\[
r^2 + \text{height}^2 = 15^2 \;\Rightarrow\; 81 + \text{height}^2 = 225
\]
\[
\text{height}^2 = 144 \;\Rightarrow\; \boxed{\text{height} = 12 \,\text{cm}}
\]

(d)(i) $r$ in terms of $h$\\[4pt]
The water surface forms a cone similar to the whole cone, so corresponding lengths are in the same ratio:
\[
\frac{r}{h} = \frac{9}{12} = \frac{3}{4} \;\Rightarrow\; \boxed{r = \frac{3}{4}h}
\]

(d)(ii) Rate at which the depth is rising when $h = 4$\\[4pt]
Volume of the water, using $V = \dfrac{1}{3}\pi r^2 h$ with $r = \dfrac{3}{4}h$ so that $V$ is a function of $h$ alone:
\[
V = \frac{1}{3}\pi\left(\frac{3}{4}h\right)^{2} \times h = \frac{1}{3}\pi \cdot \frac{9}{16}h^2 \cdot h = \frac{3\pi h^3}{16}
\]
Differentiate with respect to $h$:
\[
\frac{\dd V}{\dd h} = \frac{9\pi h^2}{16}
\]
Connect the rates by the chain rule:
\[
\frac{\dd h}{\dd t} = \frac{\dd V}{\dd t} \times \frac{\dd h}{\dd V} = \frac{\dd V}{\dd t} \div \frac{\dd V}{\dd h}
\]
When $h = 4$, $\dfrac{\dd V}{\dd h} = \dfrac{9\pi \times 16}{16} = 9\pi$, and $\dfrac{\dd V}{\dd t} = 27$:
\[
\frac{\dd h}{\dd t} = 27 \times \frac{16}{9\pi \times 16} = \frac{27}{9\pi} = \frac{3}{\pi} = 0.95492\ldots
\]
The exact value $\dfrac{3}{\pi}$ is carried to the last step, so no intermediate has been rounded.
\[
\boxed{\frac{\dd h}{\dd t} = 0.955 \,\text{cm}\,\text{s}^{-1}}
\]
\end{ansbox}

%%%%%%%%%%%%%%%%%%%
\examq{4037/21/M/J/25}{10}{Calculus for Kinematics}{11}
%%%%%%%%%%%%%%%%%%%
\begin{mstab}{q1color}
\multirow{3}{*}{10(a)(i)} & $[v =] \dfrac{t^3}{3} - 2t \ (+c)$ & \textbf{M1} & Or $\left[\dfrac{t^3}{3} - 2t\right]_{[3]}^{[4]}$ \\ \cline{2-4}
 & $[v =] \dfrac{t^3}{3} - 2t - \dfrac{10}{3}$ & \textbf{A1} & Or substituting in limits \newline $\left(\dfrac{4^3}{3} - 2 \times 4\right) - \left(\dfrac{3^3}{3} - 2 \times 3\right) - \dfrac{1}{3}$ \\ \cline{2-4}
 & $v = 10$ & \textbf{A1} & \\ \hline
\multirow{3}{*}{10(a)(ii)} & 2 correct terms in \newline $[s =] \dfrac{t^4}{12} - t^2 - \dfrac{10}{3}t \ (+d)$ & \textbf{M1} & \textbf{FT} an expression in form $at^3 + bt$ \newline\newline Or 2 correct terms in \newline $\left[\dfrac{t^4}{12} - t^2 - \dfrac{10}{3}t\right]_{[3]}^{[4]} - \dfrac{1}{4}$ \\ \cline{2-4}
 & $s = \dfrac{t^4}{12} - t^2 - \dfrac{10}{3}t + 12$ & \textbf{A1} & \textbf{A1} substituting in limits \newline $\left(\dfrac{4^4}{12} - 4^2 - \dfrac{10}{3} \times 4\right) - \left(\dfrac{3^4}{12} - 3^2 - \dfrac{10}{3} \times 3\right) - \dfrac{1}{4}$ \\ \cline{2-4}
 & $s = 4$ & \textbf{A1} & \\ \hline
\end{mstab}

\begin{mstab}{q1color}
\multirow{5}{*}{10(b)} & $[v =] 19t + \dfrac{5}{2}\mathrm{e}^{8-2t} \ (+c)$ oe & \textbf{M1} & \\ \cline{2-4}
 & $[v =] 19t + \dfrac{5}{2}\mathrm{e}^{8-2t} - \dfrac{137}{2}$ oe & \textbf{A1} & \textbf{FT} (\textit{their} $10$) $- \ 78.5$ \\ \cline{2-4}
 & $[s =] \dfrac{19t^2}{2} - \dfrac{5}{4}\mathrm{e}^{8-2t} - \dfrac{137}{2}t \ (+d)$ oe & \textbf{M1} & \textbf{dep} on previous \textbf{M1} \newline \textbf{FT} \textit{their} $-\dfrac{137}{2}$ provided not 0 or algebraic \\ \cline{2-4}
 & $[s =] \dfrac{19t^2}{2} - \dfrac{5}{4}\mathrm{e}^{8-2t} - \dfrac{137}{2}t + \dfrac{509}{4}$ oe, cao & \textbf{A1} & \\ \cline{2-4}
 & $s = 392$ or $392.2[49\ldots]$ & \textbf{A1} & \\ \hline
\end{mstab}

\begin{ansbox}{q1color}
(a)(i) Velocity when $t = 4$\\[4pt]
For $0 \le t \le 4$ the acceleration is $a = t^2 - 2$, and velocity is the integral of acceleration:
\[
v = \int (t^2 - 2)\,\dd t = \frac{t^3}{3} - 2t + c
\]
Use the condition $v = -\dfrac{1}{3}$ when $t = 3$:
\[
\frac{3^3}{3} - 2(3) + c = -\frac{1}{3} \;\Rightarrow\; 9 - 6 + c = -\frac{1}{3} \;\Rightarrow\; 3 + c = -\frac{1}{3}
\]
\[
c = -\frac{1}{3} - 3 = -\frac{10}{3}
\]
so
\[
v = \frac{t^3}{3} - 2t - \frac{10}{3}
\]
When $t = 4$:
\[
v = \frac{64}{3} - 8 - \frac{10}{3} = \frac{64 - 10}{3} - 8 = \frac{54}{3} - 8 = 18 - 8 = \boxed{v = 10 \,\text{m}\,\text{s}^{-1}}
\]

(a)(ii) Displacement from $O$ when $t = 4$\\[4pt]
Displacement is the integral of velocity:
\[
s = \int\left(\frac{t^3}{3} - 2t - \frac{10}{3}\right)\dd t = \frac{t^4}{12} - t^2 - \frac{10}{3}t + d
\]
Use $s = -\dfrac{1}{4}$ when $t = 3$:
\[
\frac{81}{12} - 9 - 10 + d = -\frac{1}{4} \;\Rightarrow\; \frac{27}{4} - 19 + d = -\frac{1}{4}
\]
\[
d = -\frac{1}{4} - \frac{27}{4} + 19 = -7 + 19 = 12
\]
so
\[
s = \frac{t^4}{12} - t^2 - \frac{10}{3}t + 12
\]
When $t = 4$:
\[
s = \frac{256}{12} - 16 - \frac{40}{3} + 12 = \frac{64}{3} - \frac{40}{3} - 4 = \frac{24}{3} - 4 = 8 - 4 = \boxed{s = 4 \,\text{m}}
\]

(b) Displacement from $O$ when $t = 10$\\[4pt]
For $4 \le t \le 10$ the acceleration is $a = 19 - 5\mathrm{e}^{8-2t}$. Since $\dfrac{\dd}{\dd t}\mathrm{e}^{8-2t} = -2\mathrm{e}^{8-2t}$, integrating gives
\[
v = \int\left(19 - 5\mathrm{e}^{8-2t}\right)\dd t = 19t + \frac{5}{2}\mathrm{e}^{8-2t} + c
\]
The velocity is continuous at $t = 4$, where part (a)(i) gives $v = 10$:
\[
19(4) + \frac{5}{2}\mathrm{e}^{0} + c = 10 \;\Rightarrow\; 76 + \frac{5}{2} + c = 10
\]
\[
c = 10 - \frac{157}{2} = -\frac{137}{2}
\]
so
\[
v = 19t + \frac{5}{2}\mathrm{e}^{8-2t} - \frac{137}{2}
\]
Integrate again for displacement:
\[
s = \int\left(19t + \frac{5}{2}\mathrm{e}^{8-2t} - \frac{137}{2}\right)\dd t = \frac{19t^2}{2} - \frac{5}{4}\mathrm{e}^{8-2t} - \frac{137}{2}t + d
\]
The displacement is continuous at $t = 4$, where part (a)(ii) gives $s = 4$:
\[
\frac{19(16)}{2} - \frac{5}{4}\mathrm{e}^{0} - \frac{137}{2}(4) + d = 4 \;\Rightarrow\; 152 - \frac{5}{4} - 274 + d = 4
\]
\[
d = 4 - 152 + \frac{5}{4} + 274 = 126 + \frac{5}{4} = \frac{509}{4}
\]
so
\[
s = \frac{19t^2}{2} - \frac{5}{4}\mathrm{e}^{8-2t} - \frac{137}{2}t + \frac{509}{4}
\]
When $t = 10$, the exponential is $\mathrm{e}^{8-20} = \mathrm{e}^{-12}$, which is very small:
\[
s = \frac{19(100)}{2} - \frac{5}{4}\mathrm{e}^{-12} - \frac{137}{2}(10) + \frac{509}{4} = 950 - 685 + 127.25 - \frac{5}{4}\mathrm{e}^{-12}
\]
\[
= 392.25 - 0.0000076\ldots = 392.24999\ldots
\]
The constants $-\dfrac{137}{2}$ and $\dfrac{509}{4}$ were carried as exact fractions rather than as rounded decimals, so nothing has been prematurely rounded. To 3 significant figures:
\[
\boxed{s = 392 \,\text{m}} \quad (392.2\ldots)
\]
\end{ansbox}

%%%%%%%%%%%%%%%%%%%
\examq{4037/21/M/J/25}{11}{Arithmetic \& Geometric Progressions}{8}
%%%%%%%%%%%%%%%%%%%
\begin{mstab}{q2color}
\multirow{7}{*}{11(a)} & $ar + ar^2 = 168$ \ or \ $ar^3 + ar^4 = 94.5$ oe soi & \textbf{B1} & \\ \cline{2-4}
 & $\dfrac{ar^3(1+r)}{ar(1+r)} = \dfrac{94.5}{168}$ oe & \textbf{M1} & Correct equation where $a$ can be directly eliminated \\ \cline{2-4}
 & $168r^2 = 94.5$ oe & \textbf{M1} & \\ \cline{2-4}
 & $r = \dfrac{3}{4}$ only nfww & \textbf{A1} & \\ \cline{2-4}
 & $\dfrac{3}{4}a + \dfrac{9}{16}a = 168$ oe & \textbf{M1} & \textbf{FT} \textit{their} $r$ \\ \cline{2-4}
 & $a = 128$ & \textbf{A1} & \\ \cline{2-4}
 & $u_6 = \dfrac{243}{8}$ only oe, isw & \textbf{B1} & \\ \hline
11(b) & $512$ & \textbf{B1} & \\ \hline
\end{mstab}

\begin{ansbox}{q2color}
(a) The 6th term\\[4pt]
With first term $a$ and common ratio $r$, the terms are $a$, $ar$, $ar^2$, $ar^3$, $ar^4$, $ar^5$, so the two given conditions are
\[
ar + ar^2 = 168 \qquad \text{and} \qquad ar^3 + ar^4 = 94.5
\]
Factorise each left-hand side and divide the second equation by the first, so that $a$ cancels directly:
\[
\frac{ar^3(1+r)}{ar(1+r)} = \frac{94.5}{168}
\]
The factor $ar(1+r)$ cancels, leaving
\[
r^2 = \frac{94.5}{168} = \frac{9}{16} \qquad \text{(equivalently } 168r^2 = 94.5\text{)}
\]
Hence $r = \pm\dfrac{3}{4}$, and since $r > 0$ is given,
\[
r = \frac{3}{4}
\]
Substitute into the first equation:
\[
a\left(\frac{3}{4}\right) + a\left(\frac{3}{4}\right)^{2} = 168 \;\Rightarrow\; \frac{3}{4}a + \frac{9}{16}a = 168
\]
\[
\frac{12a + 9a}{16} = 168 \;\Rightarrow\; \frac{21a}{16} = 168 \;\Rightarrow\; a = \frac{168 \times 16}{21} = 128
\]
The 6th term is $u_6 = ar^5$:
\[
u_6 = 128 \times \left(\frac{3}{4}\right)^{5} = 128 \times \frac{243}{1024} = \frac{243}{8}
\]
\[
\boxed{u_6 = \frac{243}{8} = 30.375}
\]

(b) Sum to infinity\\[4pt]
Since $\lvert r \rvert = \dfrac{3}{4} < 1$ the series converges, and $S_\infty = \dfrac{a}{1-r}$:
\[
S_\infty = \frac{128}{1 - \frac{3}{4}} = \frac{128}{\frac{1}{4}} = 128 \times 4 = \boxed{512}
\]
\end{ansbox}

\end{document}